\section{Conclusion}


This paper presents a Unified Dual-Mask Physical Model to address the degradation of light field (LF) depth estimation caused by the violation of the Lambertian assumption on non-Lambertian surfaces. By explicitly modeling the physical interactions of light rays, our approach integrates material-aware angular signal decomposition with a domain-balanced training strategy. This unified framework effectively mitigates the structural fractures in Epipolar Plane Images (EPIs) and establishes a rigorous physical boundary for EPI-based depth inference, yielding an overall Mean Absolute Error (MAE) of 0.133 and demonstrating robust performance in mixed urban environments with an MAE of 0.081.

The specific contributions of this work are summarized as follows:

\begin{itemize}
 \item \textbf{Material-Aware Angular Signal Decomposition:} We developed a three-layer angular signal decomposition mechanism coupled with geometric feature extraction to classify LF materials. This module explicitly separates Lambertian and non-Lambertian regions by analyzing the continuity and variance of angular signals, providing precise material priors that guide the subsequent dual-mask physical modeling and prevent erroneous depth assignments in specular or scattering areas.
 \item \textbf{Domain-Balanced Sampling and Unified Training Strategy:} We formulated a domain-balanced sampling protocol alongside a unified depth estimation training strategy tailored for multi-domain LF data. By dynamically adjusting the sampling weights across diverse datasets, this strategy alleviates domain shift and overfitting issues, which directly supports the accuracy improvement in complex mixed urban scenarios, achieving a low MAE of 0.081.
 \item \textbf{Physical Boundary Definition and Root Cause Analysis of EPI Failure:} Through extensive empirical evaluations comprising 132 experiments, we delineated the physical boundaries where the EPI assumption breaks down and conducted a comprehensive root cause analysis. We demonstrated that while EPI-based methods remain effective for Lambertian and mixed scenes, they inherently fail in non-Lambertian environments due to the physical destruction of viewpoint consistency. This analysis identifies complex texture-induced slope ambiguity and severe data scarcity as the primary bottlenecks, establishing a theoretical foundation for transitioning beyond traditional EPI paradigms.
\end{itemize}
The established physical boundaries and the proposed unified framework provide a solid foundation for advancing light field 3D perception. Subsequent investigations will focus on developing non-EPI physical models and expanding multi-domain datasets to enhance depth inference in highly complex non-Lambertian environments.


While the proposed three-layer angular signal decomposition effectively isolates material properties, its most critical system-level advantage lies in circumventing the limitations of frequency-domain analysis in sparse light fields. Traditional methods relying on the 2D Discrete Fourier Transform (2D-DFT) inherently assume dense angular sampling to satisfy the Nyquist-Shannon sampling theorem, where the maximum representable angular frequency $f_{max}$ is bounded by:
\begin{equation}
\label{eq:eq31}
f_{max} = \frac{1}{2 \Delta \theta}
\end{equation}
where $\Delta \theta$ is the angular sampling interval. In practical light field cameras (e.g., $9 \times 9$ or $7 \times 7$ angular resolutions), sparse sampling induces severe spectral aliasing, rendering high-frequency specular signatures indistinguishable from Lambertian textures. By shifting the paradigm from frequency-domain filtering to spatial-angular geometric decomposition—extracting angular gradients and coefficients of variation—our model establishes a robust material identification mechanism that is fundamentally invariant to angular undersampling. This design choice not only reduces computational overhead but also ensures that the physical boundaries of the Epipolar Plane Image (EPI) constraints are accurately mapped, making it highly suitable for resource-constrained embedded vision systems.

Furthermore, the pixel-level dual-mask routing architecture introduces nuanced behavior at the transition boundaries between Lambertian and non-Lambertian regions. In complex scenes, specular highlights often exhibit soft edges due to surface micro-facets or slight defocus, creating a mixed pixel regime where the angular cone constraint is partially violated. Unlike scene-level masking, which applies a rigid binary heuristic, our geometric pseudo-labels facilitate a localized routing mechanism. However, this exposes a subtle limitation: at extreme sub-pixel material transitions, the dual-branch network may experience gradient conflict during backpropagation, occasionally resulting in minor depth discontinuities or edge blurring. Future iterations could incorporate a boundary-aware regularization term or leverage spatial-contextual attention to explicitly smooth depth predictions along these high-gradient mask contours, thereby preserving sharp geometric structures.

Finally, while the domain-balanced sampling strategy significantly narrows the performance gap between synthetic datasets and real-world captures, an inherent domain shift persists regarding complex inter-reflections and volumetric scattering. Synthetic renderers typically approximate non-Lambertian effects using localized Bidirectional Reflectance Distribution Function (BRDF) models, which often fail to capture global illumination phenomena such as color bleeding or multi-bounce specularities. Consequently, the residual branch of our decomposition may occasionally conflate complex inter-reflections with high-frequency sensor noise. To address this physical ambiguity, future work will explore integrating differentiable physically-based rendering frameworks, such as Neural Radiance Fields (NeRF) or 3D Gaussian Splatting, into the unified training loop. By enforcing physically plausible light transport equations as an auxiliary constraint, the network could further disentangle multi-bounce light paths, ultimately advancing light field depth estimation in highly uncontrolled environments.

Despite the demonstrated efficacy of the proposed unified framework in restoring depth reliability, several physical and computational limitations remain, outlining promising directions for future research.

First, while the three-layer angular signal decomposition successfully circumvents the spectral aliasing issues of traditional 2D-DFT at $9 \times 9$ angular resolution, its robustness degrades under extremely sparse sampling (e.g., $3 \times 3$ or $2 \times 2$). In such critically sparse regimes, the residual component $R(u,v)$ in the decomposition model $I(u,v) = I_{DC} + I_{linear} + R(u,v)$ lacks sufficient degrees of freedom to reliably isolate material-specific high-frequency reflections. This inevitably leads to a drop in the classifier's accuracy. Future work could address this bottleneck by integrating spatial-contextual attention mechanisms or cross-modal sensor fusion (e.g., polarization or near-infrared cues) to compensate for the severe deficit in angular information.

Second, the current dual-mask routing architecture processes light fields on a strictly per-frame basis, which poses significant challenges for dynamic light field video systems. non-Lambertian phenomena, particularly dynamic specular highlights and caustics, shift continuously across spatial, angular, and temporal dimensions. This complex spatio-angular-temporal coupling can induce mask flickering and temporal depth artifacts in consecutive frames, degrading the performance of downstream 3D video applications. To mitigate this, future iterations will explore spatio-temporal regularization, incorporating scene flow estimation or recurrent neural networks to enforce strict temporal smoothness in both the material masks and the resulting depth maps.

Finally, while the non-Lambertian geometric branch successfully bypasses fractured Epipolar Plane Image (EPI) structures to estimate a proxy depth, it inherently struggles with the physical ambiguity of thick translucent or volumetric scattering media, where a definitive geometric surface is mathematically ill-defined. A compelling future direction is to integrate our material-aware physical priors into implicit neural representations, such as Neural Radiance Fields (NeRF) or 3D Gaussian Splatting. By utilizing the dual-mask decomposition to explicitly guide the optimization of anisotropic Bidirectional Reflectance Distribution Functions (BRDFs) and volumetric density fields, we can jointly achieve robust depth inference and high-fidelity novel view synthesis. Extending this domain-balanced training strategy to encompass large-scale, in-the-wild video light field datasets will be crucial for deploying these physical models in real-world autonomous navigation and immersive systems.